\section{Formal Grammars} 

  Once we are done scanning, we now have a list of tokens $(t_n)$. The next job is to determine if this list of tokens is valid syntax for our language, i.e. $t_1 t_2 \ldots t_n \in L$. So, the parser's first job is to determine membership of our sequence of tokens. This seems like a pretty hard problem since $L$ can be very complex. 

  Generally, a subset of $\Sigma^\ast$ doesn't give us much structure, so we would like to define some way to pick the subset $L$ out. We can do this with \textit{grammars}. 

  \begin{definition}[Formal Grammar]
    A \textbf{formal grammar} is a 4-tuple $G = (V, \Sigma, R, S)$ consisting of the following. 
    \begin{enumerate}
      \item \textit{Non-terminal Symbols}. $V$ is a finite set consisting of \textbf{non-terminal symbols}, also known as \textbf{variables}. 
      \item \textit{Terminal Symbols}. $\Sigma$ is a finite set---disjoint from $V$---consisting of \textbf{terminals}. 
      \item \textit{Production Rule}. $R$ is a relation, i.e. is a finite subset of 
        \begin{equation}
          (V \cup \Sigma)^\ast V (V \cup \Sigma)^\ast \times (V \cup \Sigma)^\ast
        \end{equation}
        where the LHS represents the product of the languages. We usually write the relation as $\alpha \to \beta$.\footnote{In math, this is usually written $\alpha R \beta$, or $\alpha \sim \beta$ if it is an equivalence relation.}

      \item \textit{Start Symbol}. $S \in V$ is the initial variable from which derivation begins. 
    \end{enumerate}
  \end{definition}

  We can think of the production rule as a set of transformations---formally called \textit{derivations}---that you can apply to a word. 

  \begin{definition}[Derivation]
    Let $G = (V, \Sigma, R, S)$ be a formal grammar. We define a binary relation $\derives$ on the set of all strings $(V \cup \Sigma)^\ast$ as follows: 
    \begin{enumerate}
      \item A string $u$ \textbf{directly derives} $v$, written $u \derives v$, if there exists strings $\phi, \psi \in (V \cup \Sigma)^\ast$ and a production rule $(\alpha \to \beta) \in R$ such that 
        \begin{equation}
          u = \phi \alpha \psi, \qquad v = \phi \beta \psi
        \end{equation}
        In other words, we say a word $w$ is derived from $v$, written $v \derives w$, if $w$ can be obtained by replacing a part of $v$ according to a rule in $R$. 

      \item We say $u$ \textbf{derives} $v$\footnote{We also say that $v$ is in the transitive closure of $u$.}, written $u \tderives v$, if $u$ can be directly derived into $v$ in $0$ or more steps. The relation $\tderives$, called the \textbf{reflexive transitive closure} of $\derives$, is defined as follows. $u \tderives v$ if 
        \begin{enumerate}
          \item $u = v$, or 
          \item There exists a finite sequence of strings $u = w_0, w_1, \ldots, w_n = v$ such that $w_i \derives w_{i+1}$ for all $0 \leq i < n$. 
        \end{enumerate}

    \end{enumerate}
  \end{definition}

  \begin{definition}[Language Derived From Grammar]
    The \textbf{language derived from $G$}, denoted $L(G)$, is the set of all terminal strings that can be derived from $S$ using the rules in $R$. It is defined: 
    \begin{equation}
      L(G) \coloneqq \{ w \in \Sigma^\ast \mid S \derives^\ast w \}
    \end{equation}
  \end{definition}

  \begin{example}
    Let $V = \{E\}$ be the non-terminals, and $\{+, -, \texttt{num}\}$ be our terminals, which come from our lexer. Then, our production rules can look something like. 
    \begin{enumerate}
      \item $E \to E + E$. 
      \item $E \to E - E$. 
      \item $E \to \texttt{num}$. 
    \end{enumerate}
    So basically, if we have some word, then we can replace any $E$ in the word by any of the three choices on the right hand side. For example, we can do the following sequence of derivations. 
    \begin{align}
      E & \derives E + E \\ 
        & \derives E - E + E \\
        & \derives E - \texttt{num} + E\\
        & \derives E + E - \texttt{num} + E 
    \end{align}
    You can keep doing this until there is no more production rules you can apply. For context free grammars, you basically do this until there are only terminals left, e.g. $\texttt{num} + \texttt{num} - \texttt{num} + \texttt{num}$. 
  \end{example}

  \begin{definition}[Transitive Closure]
  \end{definition}

  \begin{definition}[Language Generated by Grammar]
    Given grammar $G = (V, \Sigma, R, S)$, the language $L(G)$ is the set of all terminal strings that can be derived from the start symbol $S$. 
    \begin{equation}
      L(G) \coloneqq \{ w \in \Sigma^\ast \mid S \derives^\ast w \}
    \end{equation}
  \end{definition}

  Therefore, you can basically think of the ``complexity'' or size of a language $L$ as being determined by the ``complexity'' of the generating grammar $G$. Usually, the alphabet is kept fixed, and the main contributor to the complexity are the production rules $R$. Depending on how much we restrict $R$, we get different levels of languages in the \textit{Chomsky Hierarchy}. 

  \begin{definition}[Unrestricted]
    An \textbf{unrestricted language} $L$ is a language that can be generated by some grammar $G$. 
  \end{definition}

\subsection{Context-Sensitive Grammars}

  \begin{definition}[Context Sensitive]
    A grammar is \textbf{context sensitive} if 
    \begin{equation}
      R \subset \{(\alpha, \beta) \mid \alpha, \beta \in (V \cup \Sigma)^+, |\alpha| \leq |\beta|\} \cup \{(S, \epsilon) \text{ if S does not appear on any RHS} \}
    \end{equation}
    That is, the length of the string on the right must be greater than or equal to the length of the string on the left. You cannot ``delete'' symbols are you derive. A language $L$ is \textbf{context-sensitive} if $L = L(G)$ for some context-sensitive grammar $G$. 
  \end{definition}

\subsection{Context-Free Grammars}

  For theorists, context sensitive grammars are important, but for engineers, context free grammars matter much more.  

  \begin{definition}[Context Free]
    A grammar is \textbf{context free} if 
    \begin{equation}
      R \subset V \times (V \cup \Sigma)^\ast
    \end{equation}
    That is, it asserts that the left-hand side can only contain a single non-terminal. A language $L$ is \textbf{context-free} if $L = L(G)$ for some context-free grammar $G$. 
  \end{definition}

  This allows for a parse-tree structure because each variable ``branches'' into a new string independently of its neighbors. 

  \begin{definition}[Ambiguous]
    A context-free grammar is \textbf{ambiguous} if there is a string that can have more than one valid derivation tree. 
  \end{definition}

  \begin{example}[Ambiguous Context Free Grammar]
    Let $V = \{E\}$ be the non-terminals, and $\{+, -, \texttt{num}\}$ be our terminals, which come from our lexer. Then, our production rules can look something like. 
    \begin{enumerate}
      \item $E \to E + E$. 
      \item $E \to E - E$. 
      \item $E \to \texttt{num}$. 
    \end{enumerate}
    So basically, if we have some word, then we can replace any $E$ in the word by any of the three choices on the right hand side. For example, we can do the following sequence of derivations. 
    \begin{align}
      E & \derives E + E \\ 
        & \derives E - E + E \\
        & \derives E - \texttt{num} + E\\
        & \derives E + E - \texttt{num} + E 
    \end{align}
    You can keep doing this until there is no more production rules you can apply. For context free grammars, you basically do this until there are only terminals left, e.g. $\texttt{num} + \texttt{num} - \texttt{num} + \texttt{num}$. But note that this is ambiguous since there are multiple ways to derive. For example, if we wanted to get something like \texttt{num} - \texttt{num} - \texttt{num}, we can do it in either of the following ways. 
    \begin{align}
      E & \derives E - E \\ 
        & \derives E - (E - E) \\ 
        & \derives \texttt{num} - (\texttt{num} - \texttt{num}) \\
      E & \derives E - E \\ 
        & \derives (E - E) - E \\ 
        & \derives (\texttt{num} - \texttt{num}) - \texttt{num} 
    \end{align}
  \end{example}

  This is not good, so we want to modify our grammar so that it is safe from these ambiguities. 

  \begin{example}[Non-Ambiguous Context-Free Grammar]
    If we have a grammar with the following production rules. 
    \begin{enumerate}
      \item $E \to E + \texttt{num}$
      \item $E \to E - \texttt{num}$ 
      \item $E \to \texttt{num}$
    \end{enumerate}
    Then, we have no ambiguities since there is only one possible way 
    \begin{align}
      E & \derives E - \texttt{num} \\ 
        & \derives (E - \texttt{num}) - \texttt{num} \\
        & \derives (\texttt{num} - \texttt{num}) - \texttt{num} 
    \end{align}
  \end{example}

  Note that figuring out whether a grammar is ambiguous or unambiguous is a bit tricky. It is in fact undecidable. 

\subsection{Regular Grammars}

  \begin{definition}[Regular]
    A grammar $G$ is \textbf{regular} if it is either of the following: 
    \begin{enumerate}
      \item \textit{Right Linear}. 
        \begin{equation}
          R \subset V \times (\Sigma \cup \Sigma V \cup \{e\}) 
        \end{equation}
        Meaning that every rule must look like $A \to a$ or $A \to a B$. 

      \item \textit{Left Linear}. 
        \begin{equation}
          R \subset V \times (\Sigma \cup V \Sigma \cup \{\epsilon\}) 
        \end{equation}
        Meaning that every rule must look like $A \to a$ or $A \to Ba$. 
    \end{enumerate}
    A language $L$ is \textbf{regular} if $L = L(G)$ for some regular grammar $G$. 
  \end{definition}

  \begin{theorem}[Kleene's Theorem]
    A language $L$ is regular if and only if there exists a DFA $M$ such that $L(M) = L$. 
  \end{theorem}

\subsection{Parse Trees}

  You can track the derivations used to go from the base expression to the sequence of terminals. This is perfectly represented in the parse tree data structure. 

  \begin{definition}[Parse Tree]
    The \textbf{yield} of a parse tree is...
  \end{definition}

